\section{Постановка задачи}
\label{sec:Chapter1} \index{Chapter1}

\subsection{Формальная модель входа}

Рассматривается \textbf{статическая трёхмерная сцена (или крупная модель)} $S$, готовая к загрузке в ускоряющие структуры для аппаратной трассировки (TLAS и BLAS). 

На уровне данных графического движка $S$ представляет собой набор из $R$ независимых геометрий (меш-ассетов или \texttt{RElem} в терминологии DagorEngine~\cite{dagorEngine}):
\begin{equation}
    S = \bigl(V,\; \{M_k\}_{k=1}^{R}\bigr),
\end{equation}
где
\begin{itemize}
    \item $V \in \mathbb{R}^{n_v \times 3}$ -- единый массив (буфер), содержащий пространственные 3D-координаты $n_v$ уникальных вершин. Этот массив может быть общим (разделяемым) для всей сцены/ассета;
    \item $M_k \subset \{1,\dots,n_v\}^3$ -- набор треугольников $k$-того меш-ассета (геометрии/материальной подгруппы). Каждый треугольник задаётся тройкой индексов, ссылающихся на координаты вершин в общем буфере $V$;
    \item $R$ -- общее число изначальных логических кусков геометрии (например, разделённых по материалам).
\end{itemize}

В пайплайнах игровых движков графические дизайнеры зачастую разделяют сложные объекты не по пространственному признаку, а по логическому или материальному. Например, могут выделить в отдельный меш-ассет $M_k$ все деревянные конструкции сложного здания или уровня. Геометрия может существовать в нескольких уровнях детализации (LOD), при этом для формирования структуры трассировки используется геометрия наивысшего качества.

\subsection{Модель построения BLAS}

При подготовке к аппаратному рейтрейсингу графический движок подаёт каждую геометрию $M_k$ в драйвер как независимую сущность. В соответствии с этим билдер драйвера воспринимает каждый $M_k$ как единое целое и строит для него \emph{свой собственный} отдельный BLAS:
\begin{equation}
    \mathcal{B}\bigl(M_k\bigr) \;\longmapsto\; T_k,
\end{equation}
где $T_k$ -- бинарное дерево BVH, полученное из внутреннего билдера драйвера $\mathcal{B}$. 

При добавлении в сцену каждый такой построенный BLAS $T_k$ регистрируется в структуре верхнего уровня (TLAS) в виде отдельного инстанса (листового узла). Таким образом, составная сцена $S$ представлена в TLAS набором из $R$ независимых инстансов.

\subsection{Гипотеза проблемы}

Проблема возникает из-за способа группировки треугольников внутри исходных $M_k$. Поскольку каждый отдельный BLAS строится для всего подмеша $M_k$ целиком, качество построенного дерева $T_k$ и размер его ограничивающего объёма (AABB) напрямую зависят от того, насколько компактен этот подмеш в пространстве.

Если $M_k$ представляет собой логическую материальную группу (например, все деревянные конструкции здания), то его треугольники разбросаны по всему объёму здания. В результате:
1. Корневой AABB этого BLAS $T_k$ раздувается до размеров всего здания.
2. Внутри самого дерева $T_k$ получаются сильно перекрывающиеся внутренние узлы, так как SAH-билдеру сложно разделить пространственно-разреженную геометрию на неперекрывающиеся компактные объёмы.

Наличие в $S$ множества подобных подмешей (выделенных под разные материалы: металл, бетон, стекло и т.д.) приводит к тому, что в листьях TLAS образуется набор из $R$ огромных инстансов. Они полностью перекрывают друг друга в пространстве. При трассировке луч вынужден проверять пересечение практически со всеми этими инстансами, что приводит к существенному росту ложных пересечений и может существенно снижать производительность рендеринга.

\subsection{Содержательная задача}

Построить отображение $\pi$, которое преобразует исходный набор «проблемных» разбросанных геометрий $\{M_k\}$ в новый набор из $K$ пространственно-компактных мини-ассетов $\{M'_j\}$:
\begin{equation}
    \pi:\ \bigl(V,\; \{M_k\}_{k=1}^{R}\bigr) \;\longmapsto\; \bigl(\{V'_j\},\; \{M'_j\}_{j=1}^{K}\bigr).
\end{equation}
Метод извлекает все треугольники исходной геометрии (забывая об их изначальном разделении по материалам) и перегруппировывает их в $K$ новых кластеров $\{M'_j\}$ исключительно на основе их пространственной близости. Каждый кластер $M'_j$ получает собственные буферы и регистрируется в движке как совершенно новая, независимая геометрия. Драйвер, в свою очередь, построит для каждого из них свой компактный BLAS.

Набор новых ассетов должен удовлетворять следующим условиям:
\begin{enumerate}
    \item \textbf{Визуальная эквивалентность}. Для любого направления
      обзора результирующий рендер совпадает с исходным попиксельно с точностью до численных ошибок при той
      же камере и той же реализации шейдеров. Достаточное условие:
      объединение треугольников всех новых ассетов $M'_j$ совпадает
      (как мульти-множество) с множеством треугольников исходной составной сцены $S$.
    \item \textbf{Пространственная локализация}. Физический размер каждого нового мини-ассета должен быть существенно меньше размера всей исходной сцены. Формально, отношение диагонали ограничивающего параллелепипеда ($\mathrm{diag}(\mathrm{AABB})$) нового ассета к диагонали всей исходной сцены должно быть ограничено:
      \[
          \dfrac{\mathrm{diag}(\mathrm{AABB}(M'_j))}{\mathrm{diag}\left(\bigcup_k \mathrm{AABB}(M_k)\right)} \;\le\; \tau,
      \]
      где $\tau \in (0, 1)$ -- желаемый коэффициент уменьшения габаритов. Это гарантирует, что новые инстансы будут компактными и не будут сильно перекрываться друг с другом.
    \item \textbf{Контракт DXR}. Каждый новый меш-ассет $M'_j$ удовлетворяет
      ограничениям интерфейса DXR-билдера: непрерывность
      vertex/index-буферов, число вершин $\le 65{,}535$ для
      16-битного индекса (Dagor использует
      \texttt{uint16\_t}), а также лимиту на число внутренних подмешей (geometry-описаний) внутри одного BLAS: $\le M_{\max}\!=\!32$.
    \item \textbf{Снижение реального времени трассировки}. За счёт устранения перекрытий в TLAS (и внутри самих деревьев BLAS), время
      \texttt{rtsm::do\_trace} после применения метода удовлетворяет
      $t_{after} \le t_{before}\cdot(1 - \alpha)$ с $\alpha > 0$,
      демонстрирующим устойчивое снижение в условиях фиксированной камеры и
      установившегося состояния денойзера.
\end{enumerate}

\subsection{Технические требования}

\begin{description}[leftmargin=1.3cm,style=nextline]
\item[Платформа] Windows 10/11~x64, DirectX~12, DXR Tier 1.1,
  GPU NVIDIA Turing+ или AMD RDNA2+. Эталонная конфигурация
  замеров -- GeForce RTX-класса с поддержкой
  hardware-traversal.
\item[Окружение] DagorEngine, семпл \texttt{testGI}. Семпл
  поддерживает прямую загрузку моделей в \texttt{.dag}-формате
  через сценарии и обладает встроенным GPU-профайлером с
  таймингами по подсистемам.
\item[Целевая сцена] Модель \href{https://sketchfab.com/3d-models/house-ko-bulon-don-urak-lawoi-village-thailand-bc53e24d35a9439cac420cf530e7de3f}{\texttt{House, Ko Bulon Don Urak Lawoi Village, Thailand}} -- созданный
  специально для эксперимента «плохой» вариант жилого дома:
  $\sim 1{,}2\!\cdot\!10^6$ треугольников, $\sim 10$ материальных
  групп, геометрически разнесённых по всему объёму модели; в
  сцене размещены $N\!=\!16$ инстансов.
\item[Метрики измерения] (а) медианное время GPU-трассировки теней
  \texttt{rtsm::do\_trace}, измеренное через встроенный
  GPU-профайлер DagorEngine; (б) визуальная картина
  Acceleration-Structure в Nsight~Graphics с раскраской по
  индексу BLAS-инстанса для оценки качества пространственной
  локализации; (в) визуальное сравнение качества теней до и
  после оптимизации в финальном изображении.
\item[Аппаратное ограничение DXR Tier 1.1] Index-буферы BLAS в
  Dagor выделяются как \texttt{uint16\_t}, что ограничивает число
  вершин на одну BLAS-geometry значением $2^{16}-1$ и, при
  отсутствии sharing-а вершин, число треугольников значением
  $\lfloor (2^{16}-1)/3 \rfloor = 21\,845$. Метод обязан
  учитывать это ограничение.
\end{description}

\subsection{Критерии приёмки}

Работа считается выполненной успешно, если выполнено каждое из:

\begin{enumerate}
    \item Реализован метод пространственной кластеризации ассетов как
      опциональный режим инициализации BVH для динамических
      моделей в семпле \texttt{testGI}, включаемый из
      \texttt{settings.blk}; собирается штатным
      jam-скриптом DagorEngine.
    \item На модели \href{https://sketchfab.com/3d-models/house-ko-bulon-don-urak-lawoi-village-thailand-bc53e24d35a9439cac420cf530e7de3f}{\texttt{House, Ko Bulon Don Urak Lawoi Village, Thailand}} (16 инстансов в сцене)
      зафиксировано снижение медианного времени
      \texttt{rtsm::do\_trace} не менее чем в $2$ раза по
      сравнению с baseline-режимом без предобработки.
    \item Визуальная картина Acceleration-Structure в
      Nsight~Graphics с раскраской по индексу инстанса
      демонстрирует пространственно-связные области (а не
      «шумовую» картину перекрытий), подтверждая
      пространственную локализацию BLAS-кластеров.
    \item Сцена после применения метода визуально неотличима от
      сцены до его применения: тени, освещение и
      раскраска ассетов совпадают.
\end{enumerate}
